\chapter{Gravitation universelle et poids}\label{ch:g10:universal-gravitation}

Lâchez une pomme~: elle tombe. La Lune, à soixante rayons terrestres,
ne semble jamais tomber. La grande intuition de Newton est que la Lune
\emph{tombe} --- d'environ $1{,}4$~millimètre chaque seconde, exactement
la quantité qu'une attraction s'affaiblissant comme le carré de la
distance prédit. Une seule formule, avec une constante universelle,
gouverne la pomme, la Lune, les marées et les planètes. Le premier
volume a décrit la \omterm{def:g10:universal-gravitation:force}{gravitation} avec des mots~; ce chapitre la
mesure.

\section{La loi de la gravitation universelle}

\begin{definition}[Force gravitationnelle]\label{def:g10:universal-gravitation:force}
Deux corps quelconques s'attirent, quelle que soit leur composition et
quelle que soit la distance qui les sépare. Cette attraction est la
\emph{force gravitationnelle}\index{force gravitationnelle}, et le
phénomène est la \emph{gravitation}\index{gravitation}. Elle n'a besoin
d'aucun contact, d'aucune corde ni d'aucun milieu, et elle est
\emph{mutuelle}~: chaque corps tire l'autre.
\end{definition}

\begin{theorem}[Loi de la gravitation universelle]\label{thm:g10:universal-gravitation:law}
Deux corps de masses $m_1$ et $m_2$ (en \unit{kg}), dont les centres
sont séparés d'une distance $d$ (en \unit{m}), s'attirent par deux
forces dirigées le long de la droite joignant les centres, chacune
pointant vers l'autre corps, d'égale intensité
\[
F = G\,\frac{m_1 m_2}{d^2}
\qquad \text{(en \unit{N})},
\]
où $G$ est la même constante pour toute paire de corps de l'univers.
\end{theorem}
\admitted

\begin{remark}[D'où vient la loi]
Newton a distillé cette loi du mouvement observé des planètes (lois de
Kepler)~; la déduction honnête est menée dans le volume de
1\textsuperscript{re} année. Le~\cref{ch:g12:satellites-kepler}
mettra déjà la loi au travail sur les \omterm{def:g10:universal-gravitation:satellite}{satellites} et les planètes.
\end{remark}

\begin{definition}[La constante de gravitation]\label{def:g10:universal-gravitation:constant}
La constante
\[
G = \qty{6.67e-11}{N.m^2/kg^2}
\]
est la \emph{constante de gravitation}\index{constante de gravitation}.
Sa valeur minuscule est la raison pour laquelle la gravité entre
objets du quotidien est imperceptible~: ce n'est que lorsqu'au moins
une masse est astronomique que $G m_1 m_2 / d^2$ compte pour quelque
chose.
\end{definition}

\begin{omfigure}
\begin{tikzpicture}[scale=1.0]
  \draw[omExo, dashed] (0,0) -- (5.2,0);
  \fill[omDef!25] (0,0) circle (0.45);
  \draw[omDef, thick] (0,0) circle (0.45);
  \fill[omDef!25] (5.2,0) circle (0.28);
  \draw[omDef, thick] (5.2,0) circle (0.28);
  \node[font=\small] at (0,0) {$m_1$};
  \node[font=\small] at (5.2,0) {$m_2$};
  \draw[-Stealth, omThm, very thick] (0.45,0) -- (1.55,0)
    node[midway, above, font=\small] {$\vect F_{2/1}$};
  \draw[-Stealth, omThm, very thick] (4.92,0) -- (3.82,0)
    node[midway, above, font=\small] {$\vect F_{1/2}$};
  \draw[Stealth-Stealth, thin] (0,-0.75) -- (5.2,-0.75)
    node[midway, below, font=\small] {$d$};
\end{tikzpicture}

{\small Les deux forces d'une paire gravitationnelle~: le long de la
droite des centres, vers l'autre corps, avec la \emph{même} intensité
$F = G m_1 m_2 / d^2$ même lorsque les masses sont très inégales.}
\end{omfigure}

\begin{remark}[Lire la formule]\label{rem:g10:universal-gravitation:reading}
Quatre traits méritent l'attention.
\begin{itemize}
  \item \emph{Mutuelle et égale.} La Terre vous tire avec votre
        \omterm{def:g10:universal-gravitation:weight}{poids} --- et vous tirez la Terre avec exactement la
        même force. (La Terre répond simplement moins~: elle est environ
        $\num{e23}$~fois plus lourde.)
  \item \emph{Inverse du carré.} Doubler $d$ divise la force par $4$~;
        s'éloigner dix fois la divise par $100$.
  \item \emph{Distance entre centres.} Pour des corps sphériques ---
        étoiles, planètes, boulets --- $d$ se mesure de centre à centre,
        non de surface à surface.
  \item \emph{Les deux masses comptent.} Doubler l'une ou l'autre
        masse double la force~; la formule est symétrique en $m_1$ et
        $m_2$.
\end{itemize}
\end{remark}

\begin{example}[Deux amis]\label{ex:g10:universal-gravitation:friends}
Deux amis de \qty{70}{kg} se tiennent centres distants de
\qty{1.0}{m}~:
\[
F = \num{6.67e-11} \times \frac{70 \times 70}{1.0^2}
\approx \qty{3.3e-7}{N}.
\]
Le \omterm{def:g10:universal-gravitation:weight}{poids} de chaque ami est d'environ \qty{6.9e2}{N} --- deux
\emph{milliards} de fois plus grand. Leur attraction mutuelle est
réelle, mais impossible à sentir.
\end{example}

\begin{example}[La Terre retient la Lune]\label{ex:g10:universal-gravitation:earthmoon}
Avec $m_1 = \qty{5.97e24}{kg}$ (Terre), $m_2 = \qty{7.35e22}{kg}$ (Lune)
et $d = \qty{3.84e8}{m}$~:
\[
F = \num{6.67e-11} \times
\frac{\num{5.97e24} \times \num{7.35e22}}{(\num{3.84e8})^2}
\approx \qty{2.0e20}{N}.
\]
Et la Lune tire la Terre en retour avec les mêmes \qty{2.0e20}{N} ---
une traction à laquelle les océans répondent deux fois par jour, sous
forme de marées.
\end{example}

\begin{method}[Calculer une force gravitationnelle]\label{met:g10:universal-gravitation:compute}
\begin{enumerate}
  \item Convertir les masses en \omterm{def:g10:orders-of-magnitude:unit}{kilogrammes} et la distance de
        centre à centre en mètres.
  \item Appliquer $F = G m_1 m_2 / d^2$.
  \item Vérifier la puissance de dix séparément des chiffres, comme
        au~\cref{ch:g10:orders-of-magnitude}~: les \omterm{def:g10:universal-gravitation:force}{forces
        gravitationnelles} vont de $\num{e-7}$~newtons (deux personnes
        à un \omterm{def:g10:orders-of-magnitude:unit}{mètre}) à $\num{e22}$~newtons (Soleil sur Terre),
        donc un \omterm{def:g10:orders-of-magnitude:scinot}{exposant} mal placé est facile à attraper.
\end{enumerate}
\end{method}

\section{Poids et intensité de la gravité}

\begin{definition}[Poids]\label{def:g10:universal-gravitation:weight}
Le \emph{poids}\index{poids} $\vect P$ d'un objet à la surface d'un
astre (une planète, une lune, \dots) est la \omterm{def:g10:universal-gravitation:force}{force
gravitationnelle} que l'astre exerce sur lui. Il pointe vers le centre
de l'astre --- la direction que l'on appelle
\emph{verticale}\index{verticale} --- et se mesure en newtons.
\end{definition}

\begin{proposition}[Poids à la surface d'un astre]\label{prop:g10:universal-gravitation:surfaceg}
À la surface d'un astre sphérique de masse $M$ et de rayon $R$, un objet
de masse $m$ a pour \omterm{def:g10:universal-gravitation:weight}{poids}
\[
P = m\,g,
\qquad\text{où}\qquad
g = \frac{G M}{R^2}
\]
ne dépend que de l'astre, non de l'objet.
\end{proposition}

\begin{proof}
Le centre de l'objet et le centre de l'astre sont séparés de $R$, donc
la loi de la \omterm{def:g10:universal-gravitation:force}{gravitation} universelle
(\cref{thm:g10:universal-gravitation:law}) donne
$P = G M m / R^2 = m \times \left(G M / R^2\right)$. Le crochet est le
même pour tout objet à la surface~: on l'appelle $g$.
\end{proof}

\begin{definition}[Intensité du champ de gravitation]\label{def:g10:universal-gravitation:fieldstrength}
La quantité $g = G M / R^2$, en newtons par \omterm{def:g10:orders-of-magnitude:unit}{kilogramme}, est
l'\emph{intensité du champ de gravitation}\index{intensité du champ de
gravitation} à la surface de l'astre~: l'astre tire sur chaque
\omterm{def:g10:orders-of-magnitude:unit}{kilogramme} avec $g$~newtons. Sur Terre,
$g = \qty{9.81}{N/kg}$.
\end{definition}

\begin{example}[Calculer $g$]\label{ex:g10:universal-gravitation:gvalues}
Pour la Terre ($M = \qty{5.97e24}{kg}$, $R = \qty{6.37e6}{m}$)~:
\[
g = \frac{\num{6.67e-11} \times \num{5.97e24}}{(\num{6.37e6})^2}
= \frac{\num{3.98e14}}{\num{4.06e13}}
\approx \qty{9.81}{N/kg}.
\]
Pour la Lune ($M = \qty{7.35e22}{kg}$, $R = \qty{1.74e6}{m}$) le même
calcul donne $g_{\text{Moon}} \approx \qty{1.62}{N/kg}$ --- six fois
moins~: la Lune est bien plus légère, et le fait d'être plus petite ne
compense pas. Pour Mars, $g_{\text{Mars}} \approx \qty{3.73}{N/kg}$.
\end{example}

\begin{omfigure}
\begin{tikzpicture}
\begin{axis}[omaxis, axis x line=bottom, axis y line=left,
  width=8.8cm, height=5.4cm, ybar, bar width=15pt,
  symbolic x coords={Lune, Mars, Terre, Jupiter}, xtick=data,
  ylabel={$g$ (\unit{N/kg})}, ymin=0, ymax=29,
  nodes near coords,
  every node near coord/.append style={font=\small},
  enlarge x limits=0.16]
\addplot[draw=omDef, fill=omDef!25] coordinates
  {(Lune,1.62) (Mars,3.73) (Terre,9.81) (Jupiter,24.8)};
\end{axis}
\end{tikzpicture}

{\small \omterm{def:g10:universal-gravitation:fieldstrength}{Intensité du champ de gravitation} en surface sur quatre
mondes~: le même sac de sucre de \qty{1}{kg} pèse \qty{1.6}{N} sur la
Lune et \qty{25}{N} sur Jupiter.}
\end{omfigure}

\begin{omfigure}
\begin{tikzpicture}[scale=1.0]
  \fill[omExo!15] (0,0) circle (1.5);
  \draw[omExo, thick] (0,0) circle (1.5);
  \fill (0,0) circle (1.5pt) node[below, font=\small] {$O$};
  \fill (90:1.5) circle (1.8pt) node[above, font=\small] {$A$};
  \draw[-Stealth, omThm, thick] (90:1.5) -- (90:0.8)
    node[pos=0.6, right, font=\small] {$\vect P_A$};
  \fill (200:1.5) circle (1.8pt) node[left, font=\small] {$B$};
  \draw[-Stealth, omThm, thick] (200:1.5) -- (200:0.8)
    node[pos=0.6, below, font=\small] {$\vect P_B$};
  \fill (330:1.5) circle (1.8pt) node[right, font=\small] {$C$};
  \draw[-Stealth, omThm, thick] (330:1.5) -- (330:0.8)
    node[pos=0.6, above left, font=\small] {$\vect P_C$};
\end{tikzpicture}

{\small «~Le bas~» est une notion locale~: en tout point du globe le
vecteur \omterm{def:g10:universal-gravitation:weight}{poids} pointe vers le centre $O$. Deux marcheurs
antipodes se tiennent tous deux debout --- pieds l'un vers l'autre.}
\end{omfigure}

\begin{remark}[Pourquoi les astronautes rebondissent]\label{rem:g10:universal-gravitation:bounce}
Un astronaute de masse \qty{120}{kg} (corps plus combinaison) pèse
$120 \times 9.81 \approx \qty{1180}{N}$ sur Terre mais seulement
$120 \times 1.62 \approx \qty{194}{N}$ sur la Lune. Des muscles
entraînés sous \qty{1180}{N} portent soudain un sixième de la charge~:
d'où la fameuse démarche lunaire lente et bondissante. La masse --- et
avec elle l'effort nécessaire pour \emph{s'arrêter} ou \emph{tourner}
--- est inchangée, c'est pourquoi les astronautes lunaires tombaient si
souvent.
\end{remark}

\section{Masse ou poids~? La balance et le dynamomètre}

Le langage courant dit qu'un sac de farine «~pèse trois kilos~». La
physique coupe cette phrase en deux.

\begin{remark}[La masse n'est pas le poids]\label{rem:g10:universal-gravitation:massweight}
\begin{itemize}
  \item La \emph{masse} $m$ mesure la quantité de matière, en
        \omterm{def:g10:orders-of-magnitude:unit}{kilogrammes}. C'est une propriété de l'objet seul~: la
        même sur Terre, sur la Lune, ou à la dérive dans l'espace
        lointain.
  \item Le \emph{\omterm{def:g10:universal-gravitation:weight}{poids}} $P = mg$ est une force, en newtons,
        exercée \emph{par un astre} sur l'objet. Il change d'astre en
        astre avec $g$ --- et s'éteint loin de tout astre.
\end{itemize}
Les confondre est sans danger au marché et fatal dans un problème de
physique~: le \omterm{def:g10:orders-of-magnitude:unit}{kilogramme} n'est pas une \omterm{def:g10:orders-of-magnitude:unit}{unité} de force.
\end{remark}

\begin{method}[La balance et le dynamomètre]\label{met:g10:universal-gravitation:instruments}
\begin{itemize}
  \item Une \emph{balance à plateaux} compare l'objet à des masses
        étalons~: les \omterm{def:g10:universal-gravitation:weight}{poids} des deux plateaux se mettent à
        l'échelle avec le même $g$, donc la comparaison mesure la
        \emph{masse} --- même verdict sur tout astre.
  \item Un \emph{dynamomètre} (balance à ressort) mesure la force qui
        étire son ressort~: il lit le \emph{\omterm{def:g10:universal-gravitation:weight}{poids}}, en newtons, et
        sa lecture change d'astre en astre. Pour retrouver la masse,
        diviser par l'intensité locale du champ~: $m = P/g$.
\end{itemize}
\end{method}

\begin{example}[Farine sur la Lune]\label{ex:g10:universal-gravitation:flour}
Un sac de farine de \qty{3.0}{kg} s'équilibre contre \qty{3.0}{kg} de
masses étalons sur Terre \emph{et} sur la Lune. Un dynamomètre, en
revanche, lit $3.0 \times 9.81 \approx \qty{29}{N}$ sur Terre et
$3.0 \times 1.62 \approx \qty{4.9}{N}$ sur la Lune. Un marchand
vendant «~au newton~» sur la Lune donnerait six fois la farine --- le
commerce, comme la science, tourne à la masse.
\end{example}

\section{Tomber autour de la Terre~: le canon de Newton}

Lancez une pierre horizontalement~: la gravité courbe sa trajectoire en
arc et elle atterrit à quelques mètres. Lancez-la plus vite~: elle
atterrit plus loin. L'expérience de pensée de Newton pousse l'idée à sa
limite. La Terre est ronde, et sa surface s'abaisse d'environ
\qty{5}{m} sous l'horizontale pour chaque \qty{8}{km} parcourus~; un
corps en chute libre tombe de \qty{4.9}{m} dans sa première seconde
(le~\cref{ch:g12:projectile-motion} l'établira). Donc un projectile
rasant la Terre à \qty{8}{km/s} tombe exactement aussi vite que le sol
s'éloigne en se courbant~: il tombe pour toujours sans jamais atterrir.

\begin{definition}[Satellite et orbite]\label{def:g10:universal-gravitation:satellite}
Un \emph{satellite}\index{satellite} d'un astre est un corps qui tombe
autour de l'astre sans jamais atteindre sa surface~; le trajet qu'il
répète est son \emph{orbite}\index{orbite}.
\end{definition}

\begin{omfigure}
\begin{tikzpicture}[scale=1.0]
  \fill[omExo!15] (0,0) circle (1.5);
  \draw[omExo, thick] (0,0) circle (1.5);
  \fill[omExo!60] (-0.22,1.44) -- (0,1.85) -- (0.22,1.44) -- cycle;
  \coordinate (L) at (0,1.85);
  \draw[omThm, thick] (L) .. controls (0.6,1.85) .. (55:1.5);
  \draw[omProp, thick] (L) .. controls (1.1,1.85) and (1.7,1.1) .. (10:1.5);
  \draw[omMeth, thick] (L) .. controls (1.8,1.8) and (2.3,-0.4) .. (-70:1.5);
  \draw[-Stealth, omDef, thick] (0,1.85) arc (90:-250:1.85);
  \node[font=\small, omDef, anchor=east] at (-1.5,1.45) {orbite};
\end{tikzpicture}

{\small Le canon de Newton~: tiré de plus en plus vite du sommet d'une
montagne, le boulet atterrit de plus en plus loin --- jusqu'à ce qu'à
environ \qty{7.9}{km/s}, le sol s'éloigne aussi vite que le boulet
tombe, et il \omterm{def:g10:universal-gravitation:satellite}{orbite}.}
\end{omfigure}

\begin{remark}[La Lune est une pierre qui tombe]\label{rem:g10:universal-gravitation:fallingmoon}
La Lune est le \omterm{def:g10:universal-gravitation:satellite}{satellite} naturel de la Terre~: elle tombe vers
nous d'environ \qty{1.4}{mm} chaque seconde, et son mouvement latéral
l'emporte autour avant qu'elle puisse se rapprocher
(l'\cref{exo:g10:universal-gravitation:15} vous laisse refaire le propre
contrôle de Newton de ce nombre). Les astronautes d'une station en
orbite sont dans la même chute libre que leur vaisseau --- c'est cela,
et non une absence de gravité, qui les fait flotter. La
\omterm{def:g10:universal-gravitation:force}{gravitation} façonne toute trajectoire, de l'arc de la pierre
à l'ellipse de la planète~; l'histoire quantitative des
\omterm{def:g10:universal-gravitation:satellite}{orbites} est racontée au~\cref{ch:g12:satellites-kepler}.
\end{remark}

\begin{notation}[Fiche de données]
Sauf indication contraire d'un exercice, utiliser
$G = \qty{6.67e-11}{N.m^2/kg^2}$ et~:
\begin{center}
\small
\begin{tabular}{lccc}
 & masse (\unit{kg}) & rayon (\unit{m}) & $g$ (\unit{N/kg}) \\
Terre & \num{5.97e24} & \num{6.37e6} & 9.81 \\
Lune & \num{7.35e22} & \num{1.74e6} & 1.62 \\
Mars & \num{6.42e23} & \num{3.39e6} & 3.73 \\
Soleil & \num{1.99e30} & --- & --- \\
\end{tabular}
\end{center}
Distance Terre--Lune~: \qty{3.84e8}{m}~; distance Terre--Soleil~:
\qty{1.496e11}{m} (de centre à centre).
\end{notation}

\section{Exercices}

\begin{exercise}[$\star$]\label{exo:g10:universal-gravitation:1}
Deux danseurs de \qty{60}{kg} se tiennent centres distants de
\qty{0.80}{m}.
\begin{enumerate}
  \item Calculer leur attraction gravitationnelle mutuelle.
  \item La comparer au \omterm{def:g10:universal-gravitation:weight}{poids} d'un moustique de \qty{2.5}{mg}.
\end{enumerate}
\end{exercise}

\begin{exercise}[$\star$]\label{exo:g10:universal-gravitation:2}
Une élève a une masse de \qty{55}{kg}.
\begin{enumerate}
  \item Calculer son \omterm{def:g10:universal-gravitation:weight}{poids} sur Terre.
  \item Elle voyage sur la Lune. Quelles sont sa masse et son
        \omterm{def:g10:universal-gravitation:weight}{poids} là-bas~?
\end{enumerate}
\end{exercise}

\begin{exercise}[$\star$]\label{exo:g10:universal-gravitation:3}
\begin{enumerate}
  \item À partir de la fiche de données, recalculer l'\omterm{def:g10:universal-gravitation:fieldstrength}{intensité du
        champ de gravitation} à la surface de Mars.
  \item Un rover de \qty{900}{kg} y est envoyé. Calculer son
        \omterm{def:g10:universal-gravitation:weight}{poids} sur Mars et son \omterm{def:g10:universal-gravitation:weight}{poids} sur Terre.
\end{enumerate}
\end{exercise}

\begin{exercise}[$\star$]\label{exo:g10:universal-gravitation:4}
En utilisant la fiche de données, calculer la force exercée par la
Terre sur la Lune. Quelle force la Lune exerce-t-elle sur la Terre~?
\end{exercise}

\begin{exercise}[$\star$]\label{exo:g10:universal-gravitation:5}
Deux corps s'attirent avec une force $F_0$. Que devient la force si~:
\begin{enumerate}
  \item la distance entre centres est doublée~;
  \item la distance est divisée par deux~;
  \item les deux masses sont doublées (même distance)~;
  \item une masse est triplée \emph{et} la distance est triplée~?
\end{enumerate}
\end{exercise}

\begin{exercise}[$\star\star$]\label{exo:g10:universal-gravitation:6}
Deux superpétroliers chargés de \qty{5.0e8}{kg} chacun sont amarrés
centres distants de \qty{100}{m}.
\begin{enumerate}
  \item Calculer leur attraction gravitationnelle mutuelle.
  \item Quelle masse a un \omterm{def:g10:universal-gravitation:weight}{poids} égal à cette force sur Terre~?
  \item La force est étonnamment grande --- pourtant les pétroliers ne
        montrent aucune inclination à dériver l'un vers l'autre.
        Calculer l'accélération $a = F/m$ qu'elle donne à un pétrolier
        et commenter.
\end{enumerate}
\end{exercise}

\begin{exercise}[$\star\star$]\label{exo:g10:universal-gravitation:7}
Une boutique lunaire vend un sac de riz de \qty{3.0}{kg}.
\begin{enumerate}
  \item Que indiquent une balance à plateaux et un dynamomètre pour ce
        sac sur Terre~? Sur la Lune~?
  \item Un client paie pour «~\qty{3.0}{kg}~». Quel instrument garantit
        un marché équitable sur tout monde, et pourquoi~?
\end{enumerate}
\end{exercise}

\begin{exercise}[$\star\star$]\label{exo:g10:universal-gravitation:8}
La Station spatiale internationale vole à une altitude de
\qty{400}{km}.
\begin{enumerate}
  \item Calculer l'\omterm{def:g10:universal-gravitation:fieldstrength}{intensité du champ de gravitation} de la
        Terre à cette altitude. (Attention~: quelle est la distance au
        \emph{centre} de la Terre~?)
  \item L'exprimer en pourcentage de la valeur en surface.
  \item Les astronautes à bord flottent. Concilier cela avec votre
        réponse.
\end{enumerate}
\end{exercise}

\begin{exercise}[$\star\star$]\label{exo:g10:universal-gravitation:9}
À la surface de Vénus, $g = \qty{8.87}{N/kg}$, et le rayon de la
planète est \qty{6.05e6}{m}. En déduire la masse de Vénus et la
comparer à celle de la Terre.
\end{exercise}

\begin{exercise}[$\star\star$]\label{exo:g10:universal-gravitation:10}
\begin{enumerate}
  \item Calculer la force exercée par le Soleil sur \qty{1.0}{kg} d'eau
        de mer, puis la force exercée par la Lune sur le même
        \omterm{def:g10:orders-of-magnitude:unit}{kilogramme}.
  \item Le Soleil tire plus fort --- d'un facteur combien~?
  \item Pourtant la Lune domine les marées. La marée est entraînée non
        par la traction elle-même mais par la \emph{différence} entre
        les tractions sur les faces proche et lointaine de la Terre.
        Expliquer, sans calculer, pourquoi la Lune proche peut battre
        le Soleil distant sur ce critère.
\end{enumerate}
\end{exercise}

\begin{exercise}[$\star\star$]\label{exo:g10:universal-gravitation:11}
\begin{enumerate}
  \item Une planète a deux fois la masse de la Terre et deux fois son
        rayon. Exprimer son intensité de champ en surface en termes du
        $g$ terrestre, puis en \unit{N/kg}.
  \item Quel rayon faudrait-il à une planète de la masse de la Terre
        pour que son intensité de champ en surface soit $2g$~?
\end{enumerate}
\end{exercise}

\begin{exercise}[$\star\star\star$]\label{exo:g10:universal-gravitation:12}
Quelque part sur la droite Terre--Lune, les deux tractions sur une
sonde spatiale s'annulent. Soit $x$ la distance du centre de la Terre à
ce point et $D$ la distance Terre--Lune.
\begin{enumerate}
  \item Montrer qu'en ce point
        $\left(\dfrac{x}{D - x}\right)^{2}
        = \dfrac{M_{\text{Earth}}}{M_{\text{Moon}}}$.
  \item En déduire $x$. Quelle fraction du trajet vers la Lune
        est-ce~?
\end{enumerate}
\end{exercise}

\begin{exercise}[$\star\star\star$]\label{exo:g10:universal-gravitation:13}
Le canon de Newton, avec des chiffres. Un projectile rase la Terre
horizontalement à la vitesse $v$.
\begin{enumerate}
  \item Après $x = \qty{8.0}{km}$ de parcours horizontal, la Terre
        sphérique «~s'est abaissée~» d'une hauteur $h$ satisfaisant
        $(R + h)^2 = R^2 + x^2$. En utilisant
        $(R + h)^2 \approx R^2 + 2Rh$ pour $h$ petit, montrer que
        $h \approx x^2 / (2R)$ et calculer $h$.
  \item Un corps en chute libre tombe de \qty{4.9}{m} dans sa première
        seconde. Comparer à $h$ et en déduire la vitesse à laquelle le
        projectile n'atterrit jamais.
  \item Calculer la durée d'une \omterm{def:g10:universal-gravitation:satellite}{orbite} complète à cette
        vitesse, en rasant la surface.
\end{enumerate}
\end{exercise}

\begin{exercise}[$\star\star\star$]\label{exo:g10:universal-gravitation:14}
Une étoile à neutrons compacte $M = \qty{2.8e30}{kg}$ ($1{,}4$~fois la
masse du Soleil) dans un rayon de \qty{12}{km}.
\begin{enumerate}
  \item Calculer l'\omterm{def:g10:universal-gravitation:fieldstrength}{intensité du champ de gravitation} à sa
        surface.
  \item Que pèserait un humain de \qty{70}{kg} là-bas~? Comparer à son
        \omterm{def:g10:universal-gravitation:weight}{poids} sur Terre.
  \item Quel est le \omterm{def:g10:universal-gravitation:weight}{poids}, là-bas, d'un grain de sable de
        \qty{1.0}{mg}~? Quelle masse a ce \omterm{def:g10:universal-gravitation:weight}{poids} sur Terre~?
\end{enumerate}
\end{exercise}

\begin{exercise}[$\star\star\star$]\label{exo:g10:universal-gravitation:15}
Le test lunaire de Newton (1666). Le rayon de l'\omterm{def:g10:universal-gravitation:satellite}{orbite} de la
Lune est d'environ $60$ rayons terrestres.
\begin{enumerate}
  \item D'un facteur combien la traction de la Terre par
        \omterm{def:g10:orders-of-magnitude:unit}{kilogramme} est-elle plus faible à la Lune qu'à la surface
        de la Terre, si la gravité suit l'inverse du carré~?
  \item Près de la surface, un corps qui tombe chute de \qty{4.9}{m}
        dans sa première seconde. De combien la Lune devrait-elle
        tomber vers la Terre chaque seconde~?
  \item Vérifier contre l'\omterm{def:g10:universal-gravitation:satellite}{orbite}~: la Lune parcourt son cercle
        de rayon $d = \qty{3.84e8}{m}$ en $T = 27{,}3$~jours. Calculer
        sa vitesse $v$, la distance $x$ qu'elle couvre en une seconde,
        puis la chute $h \approx x^2/(2d)$ (comme dans
        l'\cref{exo:g10:universal-gravitation:13}).
  \item Comparer les réponses des questions~2 et~3, et énoncer ce que
        Newton en conclut.
\end{enumerate}
\end{exercise}

\section{Problème~: Peser la Terre}

\begin{problem}[{Problème du week-end --- de deux sphères de plomb dans
un hangar à la masse de la Terre puis du Soleil~: ce qu'une petite
constante, mesurée une fois, vaut}]
\label{pb:g10:universal-gravitation:1}
En 1798 Henry Cavendish suspendit une tige légère portant deux petites
sphères de plomb à un fil fin, rapprocha deux grosses sphères de plomb,
et mesura la torsion presque inexistante du fil. Les journaux
dirent qu'il avait \emph{pesé la Terre} --- et ils avaient raison~: une
fois $G$ connue, le $g$ et le rayon propres de la Terre livrent sa
masse, et l'\omterm{def:g10:universal-gravitation:satellite}{orbite} annuelle de la Terre livre celle du Soleil.
Ce problème retrace toute la chaîne. Utiliser la fiche de données tout
au long.

\medskip
\noindent\textbf{Partie I --- La force la plus faible.}
\begin{enumerate}
  \item Le fer de la tour Eiffel a une masse d'environ \qty{7.3e6}{kg}.
        Un touriste de \qty{55}{kg} se tient à \qty{100}{m} de son
        centre de masse. Calculer la traction de la tour sur le
        touriste.
  \item La comparer au \omterm{def:g10:universal-gravitation:weight}{poids} d'un grain de sable de
        \qty{1.0}{mg}.
  \item Exercice de proportionnalité~: que devient une
        \omterm{def:g10:universal-gravitation:force}{force gravitationnelle} si (a) la distance est multipliée
        par $10$~; (b) les deux masses sont multipliées par $1000$~;
        (c) les deux masses \emph{et} la distance sont multipliées par
        $1000$~?
  \item Calculer la traction de la Terre sur un objet de \qty{1.00}{kg}
        à sa surface, à partir de $M$, $R$ et $G$. Quel nom et quel
        symbole du quotidien ce nombre porte-t-il~?
  \item La gravité est de loin l'interaction la plus faible que vous
        ayez rencontrée --- et pourtant elle gouverne l'univers.
        Expliquer en une ou deux phrases pourquoi la traction de la
        Terre sur vous domine celle de votre voisin.
\end{enumerate}

\medskip
\noindent\textbf{Partie II --- La balance de Cavendish.}
Dans la balance de torsion, une grosse sphère de masse
$m_1 = \qty{158}{kg}$ attire une petite de masse
$m_2 = \qty{0.73}{kg}$~; leurs centres sont distants de
$d = \qty{0.23}{m}$.
\begin{enumerate}[resume]
  \item Avec la valeur moderne de $G$, calculer la force entre les
        deux sphères.
  \item Calculer le \omterm{def:g10:universal-gravitation:weight}{poids} de la petite sphère et le rapport des
        deux forces. Pourquoi Cavendish avait-il besoin d'un fil de
        torsion plutôt que d'une balance ordinaire~?
  \item La torsion mesurée par Cavendish correspond (en
        \omterm{def:g10:orders-of-magnitude:unit}{unités} modernes) à $F = \qty{1.47e-7}{N}$. En déduire
        sa valeur de $G$ et la comparer à celle d'aujourd'hui.
  \item Avec $G = \qty{6.67e-11}{N.m^2/kg^2}$, $g = \qty{9.81}{N/kg}$
        et $R = \qty{6.37e6}{m}$, en déduire la masse de la Terre.
        (C'est l'étape que les journaux célébrèrent.)
  \item En déduire la masse volumique moyenne de la Terre
        $\rho = M / V$, avec $V = \frac43 \pi R^3$. Les roches de
        surface ont des masses volumiques autour de
        \qty{2.7e3}{kg/m^3}~: que révèle la comparaison sur
        l'intérieur profond~?
\end{enumerate}

\medskip
\noindent\textbf{Partie III --- Autres mondes.}
\begin{enumerate}[resume]
  \item À partir de la fiche de données, recalculer l'intensité du
        champ en surface de la Lune.
  \item Sur Terre, un bon élan soulève votre centre de masse de
        \qty{0.50}{m}. Pour le même élan, la hauteur atteinte est
        inversement proportionnelle à $g$ (comme le chapitre sur
        l'énergie, le~\cref{ch:g10:energy-conservation}, le
        justifiera). Quelle hauteur le même saut vous emporte-t-il sur
        la Lune~?
  \item Mercure~: $M = \qty{3.30e23}{kg}$, $R = \qty{2.44e6}{m}$.
        Calculer son intensité de champ en surface et comparer à Mars.
        Mercure est bien plus petite que Mars --- comment les deux
        valeurs peuvent-elles être si proches~?
  \item À quelle altitude au-dessus de la surface de la Terre
        l'intensité du champ est-elle tombée à la moitié de sa valeur
        en surface~?
  \item À quelle distance du centre de la Terre la traction de la
        Terre par \omterm{def:g10:orders-of-magnitude:unit}{kilogramme} tombe-t-elle à la valeur en surface
        de la Lune, \qty{1.62}{N/kg}~? L'exprimer en rayons
        terrestres.
\end{enumerate}

\medskip
\noindent\textbf{Partie IV --- Peser le Soleil.}
La Terre parcourt une \omterm{def:g10:universal-gravitation:satellite}{orbite} quasi circulaire de rayon
$d = \qty{1.496e11}{m}$ en une année, $T = \qty{3.156e7}{s}$. On
admet (cela se prouve au~\cref{ch:g12:satellites-kepler}) qu'un corps
sur un trajet circulaire de rayon $d$ à la vitesse $v$ a besoin d'une
traction de $v^2/d$ newtons par \omterm{def:g10:orders-of-magnitude:unit}{kilogramme} vers le centre.
\begin{enumerate}[resume]
  \item Calculer la vitesse orbitale $v$ de la Terre.
  \item En déduire la traction par \omterm{def:g10:orders-of-magnitude:unit}{kilogramme} que le Soleil
        exerce sur la Terre.
  \item Cette traction est aussi $G M_{\text{Sun}} / d^2$. En déduire
        la masse du Soleil.
  \item Combien de Terres est-ce~? Comparer $M_{\text{Sun}}$ à
        $M_{\text{Earth}}$.
  \item Final. Les sphères de Cavendish, la Terre, le Soleil~: une
        seule constante $G$ a servi pour les trois. Énoncer en une
        phrase ce que le mot \emph{universelle} dans «~\omterm{def:g10:universal-gravitation:force}{gravitation}
        universelle~» nous a acheté dans ce problème.
\end{enumerate}
\end{problem}
